\section{\rev{평가 설계}}
\label{sec:experiments}

\begin{revision}

\subsection{기존 회귀시험과 자연 검토}
\label{sec:legacy-test}
기존 40개의 기준--변형 쌍은 각 규칙에 10쌍씩 배정하고 기대 판정이 나오는 쌍만 남긴 규칙 실행 회귀시험이다. P1--P3의 주석/구조화 사건 시험 30쌍과 P4 처리 시험 10쌍을 분리한다. 이 회귀 결과를 자연 오류 precision/recall/F1로 제시하지 않는다. 이전 요구를 제거한 현재 사건 검사는 기억 효과를 확인하는 ablation이다.

자연 검토는 선정된 27장면에서 네 검토자의 시간 요구ㆍ행동열ㆍ텍스트 일관성 범주를 재집계한다. 원 라벨은 변경하지 않고 Fleiss' $\kappa$와 평균 쌍별 일치율, 검토자별 분포를 보고한다. 다수가 선택한 라벨이나 합의 CSV를 새 외부 정답으로 삼지 않는다. 이번 개정에서 사람 재평가ㆍ독립 중재와 자연 오류 정답 수집을 추가하지 않았다.

\subsection{144개 통제 시나리오의 최초 실행}
\label{sec:scenario-design}
6상황의 8변형을 중간 사건 0/2/6개와 교차시켜 144개를 모두 실행했다. 각 변형은 18건이고 기대 분포는 모순 72, 양립 36, 근거 부족 36이다. 기대 판정ㆍ최초 사건ㆍ의무ㆍ사유는 최초 실행 전에 고정했으며 결과에 따른 선별은 하지 않았다. 원문과 프레임은 동시 작성했으므로 자연 문장 자동 추출 시험과 구분한다. 첫 결과에 맞지 않은 48개도 모두 보존했다.

\begin{table}[t]
\centering
\color{revisionblue}
\bitablecaption{시나리오 변형과 고정된 기대 판정(변형별 18건).}{Scenario variants and fixed expected verdicts, 18 cases per variant.}
\label{tab:scenario-variants}
\scriptsize
\begin{tabularx}{\columnwidth}{L{.45\columnwidth}Y}
\toprule
변형 & 기대 판정의 근거 \\
\midrule
현재 미해제 후 진행 & A=F, GO: 모순 \\
이전 구간 미해제 후 진행 & A=F의 유효 구간 유지, GO: 모순 \\
정상 해제 & A=T 후 GO: 양립 \\
해제 증거 누락 & A=?, B=T, GO: 근거 부족 \\
다른 대상 해제 & OTHER=T, A=?, GO: 근거 부족 \\
행동 순서 위반 & A=T이나 선행 확인=F, GO: 모순 \\
복수 의무 일부 해제 & A=T, C=F, GO: 모순 \\
복수 의무 모두 해제 & A=T, C=T, GO: 양립 \\
\bottomrule
\end{tabularx}
\end{table}

비교 경로는 동결된 기존 Python 전체 이력 검사, 현재 사건 검사, 이전 1/3/5개 사건 창, 별도 FSM, 기존 UPPAAL, 기존 자동 텍스트 경로다. 창 비교는 기존 검사기에 잘린 이력을 제공한 ablation이며 독립 경쟁 알고리즘은 아니다. 구조화 기존 경로에는 단일 의무 투영, FSM에는 조건별 프레임을 제공한다. 따라서 그 차이에는 기억 길이뿐 아니라 의무 표현 용량과 미상 누적 정책도 포함된다. 단일 의무를 지원하는 경로의 복수 의무 도전 사례도 분모에서 제외하지 않는다.

이 설계에서 간격은 초 단위 시간이 아니라 사이에 끼운 사건의 수다. 현재 미해제 변형은 진행 사건에 A=F가 다시 명시되지만, 이전 미해제 변형은 앞서 확인된 A=F를 중간 사건 뒤까지 보존해야 한다. 따라서 같은 GO라도 현재 정보로 충분한 사례와 이전 증거가 필요한 사례를 나눠 볼 수 있다. ``다른 대상 해제''는 다른 조건의 참 증거가 A를 대신 해제하지 않는지, 복수 일부/모두 해제는 하나의 참 증거가 다른 의무를 지우거나 남기지 않는지 확인한다. 정상ㆍ근거 부족 변형도 함께 두어 확정 모순만 찾는 시험으로 축소하지 않았다.

FSM은 active 의무 집합과 조건별 알려진 증거를 갱신하는 별도 단순 구현이다. 작성된 조건 상태를 입력받으므로 자연어 의미를 독립 판정하는 oracle이 아니다. 기존 자동 경로에는 원문만 제공하고 증거 정답을 넣지 않았다. 세 분류 일치와 모순 72건을 분모로 한 최초 위반 위치 일치를 보고하며, 모순 사례를 미상으로 남긴 것도 실패 분모에 포함한다.

세 분류 일치는 기대 판정과 최종 판정이 같은 사례 수다. 최초 위반은 기대 모순72건에서 정확한 사건 ID를 돌려준 수이므로 분류를 맞혔더라도 위치가 다르면 일치로 세지 않는다. 문제 위치108건은 모순72와 근거 부족36을 합친 별도 추적 대상이다. 관련 의무와 사유도 각각 대조하며, 자연 자료의 라벨과 이 작성 기대 기록을 교환하지 않는다. 이 구분은 어떤 경로가 문제를 판정하는지와 수정에 필요한 위치 정보를 실제로 돌려주는지를 나눠 확인하게 한다.

\subsection{동일 사례 재시험과 구현 일치성}
\label{sec:retest-design}
48개 최초 불일치의 원인을 분석한 뒤 조건별 UPPAAL 모델을 개발했다. 기존 코드와 기대 판정은 바꾸지 않고 같은 \texttt{cases.json}을 재사용했다. 개발 후 동일 사례 재시험이며 개발 미사용 독립 시험이 아니다. 모순 72건의 최초 위반과 모순/근거 부족 108건의 문제 위치ㆍ관련 의무ㆍ사유를 비교했다. 성공만 보고하지 않도록 최초 시험과 재시험을 별도 표에 둔다.

기존 Python--UPPAAL 비교는 판정ㆍ오류 종류ㆍ최초 위치ㆍ미상 사유를 항목별로 대조한다. 사유가 다른 사례는 최초 반례와 종료 상태를 따로 읽어 모델 의미 문제와 출력 범위 문제를 구분한다. 전체 상태 경로의 동등성이나 규칙 자체의 독립 타당성을 검증한 시험은 아니다.

\subsection{보조 분석과 통계의 범위}
\label{sec:aux-design}
속도 반응 분석은 $W=0.5$초 구간에서 속도 변화량 $\delta=0.1$~m/s 이상이고 방향 일관 비율 $\rho=0.7$ 이상인 시작점을 찾으며 $D=3$초 제한과 비교한다. 정지ㆍ양보의 $v\leq0.5$~m/s 이미 만족 규칙은 초기 후보 검토 뒤 추가되었다. 이 값들은 법규ㆍ차량 안전 기준에서 도출한 것이 아니라 탐색 설정이다. 기존 27설정의 130건을 항상 검출ㆍ설정 의존ㆍ항상 미검출로 재집계했다. 반복 CoC의 최초 시각 보존과 재시작, 장면 경계 연결 처리도 별도 보조 결과로 유지한다.

144건은 동일 문장 틀과 변형을 교차한 결정적 조합 시험이다. 자연 모집단의 무작위 표본이나 독립 시행이 아니므로 모집단 오류율과 95\% 신뢰구간을 산출하지 않는다. 자연 검토도 확정된 독립 오류 정답이 없어 TP/FP/TN/FN, sensitivity, specificity, precision/recall/F1을 계산하지 않는다. 미상을 정상으로 합치지 않고 실제 분류 건수를 함께 보고한다.
\end{revision}
